% Canonical terminology. One macro per defined object.
%
% Prose uses the macro, never a hand-typed synonym. This enforces the banana
% rule mechanically: scripts/check_prose_quality.py reports any hand-typed
% occurrence of a term that has a macro here.
%
% Rename a concept in exactly one place: here. Every default below is a
% placeholder the writer replaces; a macro body must be a distinctive
% technical object (a method, a factor, a metric, a dataset), never a
% generic phrase such as "the model".

% --- study identity -------------------------------------------------------
\newcommand{\projectname}{\slot{Project name}\xspace}
\newcommand{\methodname}{\slot{Method name}\xspace}

% --- compared arms --------------------------------------------------------
\newcommand{\armprimary}{\slot{primary arm}\xspace}
\newcommand{\armbaseline}{\slot{strongest baseline}\xspace}
\newcommand{\armcontrol}{\slot{control or second baseline}\xspace}

% --- experimental factors ------------------------------------------------
% PLURALS: every macro here ends in \xspace, which inserts a space before a
% following letter. `\condition s` therefore renders as "dataset s", and
% `\condition{}s` does the same. It has shipped in a draft and no gate but
% a reader caught it. Write the plural as literal text, or give it its own
% macro. check_prose_quality.py now refuses `\macro s`.
\newcommand{\condition}{\slot{condition factor, e.g. data split}\xspace}
\newcommand{\repetition}{\slot{repetition factor, e.g. training seed}\xspace}

% --- quantities ----------------------------------------------------------
\newcommand{\metricname}{\slot{primary metric}\xspace}

% --- symbols -------------------------------------------------------------
% Defined once here, used identically in prose, equations, tables, captions.
\newcommand{\dataset}{\mathcal{D}}         % data set
\newcommand{\model}{f_\theta}              % the model

% --- compatibility aliases for the legacy page-1 card (figures/hero.tex) ----
% Only figures/hero.tex (the never-drawn design-card fallback) still uses
% these names. They resolve to the generic arms and factors above.
\newcommand{\methodcontrol}{\armcontrol}
\newcommand{\methodmax}{\armbaseline}
\newcommand{\methodrandom}{\armprimary}
\newcommand{\splitseed}{\condition}
\newcommand{\trajseed}{\repetition}
\newcommand{\evalsem}{evaluation semantics\xspace}
\newcommand{\certbound}{\metricname}
